% !TEX root = frode-bibelen.tex
% ----------------------------------------------------------------------------
%  Frode-bibelen -- oppsett: pakker, farger, kodestil, bokser og topptekst
%  Inkluderes fra frode-bibelen.tex (kompiler alltid hovedfilen).
% ----------------------------------------------------------------------------

\usepackage[utf8]{inputenc}
\usepackage[T1]{fontenc}
\IfFileExists{norsk.ldf}
  {\usepackage[norsk,shorthands=off]{babel}}
  {\usepackage[english,shorthands=off]{babel}
   \addto\captionsenglish{\renewcommand{\contentsname}{Innhold}%
     \renewcommand{\figurename}{Figur}\renewcommand{\tablename}{Tabell}}}

\usepackage{charter}                       % brødtekst
\usepackage[scaled=0.92]{helvet}           % overskrifter
\renewcommand{\ttdefault}{txtt}            % kode
\usepackage{microtype}
\usepackage{pifont}

\usepackage[a4paper,margin=1.8cm,top=2.1cm,bottom=2.0cm,headsep=0.5cm]{geometry}
\usepackage[table]{xcolor}
\usepackage{listings}
\usepackage[skins,breakable]{tcolorbox}
\usepackage{fancyhdr}
\usepackage{titlesec}
\usepackage{enumitem}
\usepackage{booktabs}
\usepackage{array}
\usepackage{longtable}
\usepackage{multicol}
\usepackage{tikz}
\usetikzlibrary{positioning,arrows.meta,calc,shapes.multipart}
\usepackage[hidelinks]{hyperref}

% ---------------------------------------------------------------- farger ---
\definecolor{frodeblue}{HTML}{1F4E79}
\definecolor{kw}{HTML}{1F3F9F}
\definecolor{cm}{HTML}{4F7A4F}
\definecolor{st}{HTML}{9A2A2A}
\definecolor{bg}{HTML}{F5F7FA}
\definecolor{skinn}{HTML}{1C120E}          % mørkt skinn (forside)
\definecolor{skinnlys}{HTML}{3B271D}
\definecolor{gull}{HTML}{D4AF37}
\definecolor{gullmork}{HTML}{A88A2E}

% Forside: \morkforsidefalse = hvit bakgrunn med gullramme (for utskrift).
%          \morkforsidetrue  = svart skinn og gull (for skjerm/PDF).
\newif\ifmorkforside
\morkforsidefalse

% ------------------------------------------------------------- listings ---
\lstdefinestyle{frode}{
  language=C,
  basicstyle=\ttfamily\footnotesize,
  keywordstyle=\bfseries\color{kw},
  commentstyle=\color{cm},
  stringstyle=\color{st},
  morekeywords={bool,true,false,NULL,FILE},
  showstringspaces=false,
  columns=fullflexible, keepspaces=true,
  tabsize=2,
  breaklines=true,
  backgroundcolor=\color{bg},
  frame=leftline, framerule=2pt, rulecolor=\color{frodeblue!45},
  xleftmargin=8pt, framexleftmargin=5pt,
  aboveskip=6pt, belowskip=6pt,
  upquote=true,
  literate={æ}{{\ae}}1 {ø}{{\o}}1 {å}{{\aa}}1
           {Æ}{{\AE}}1 {Ø}{{\O}}1 {Å}{{\AA}}1
           {é}{{\'e}}1 {É}{{\'E}}1 {–}{{--}}1
           {«}{{\guillemotleft}}1 {»}{{\guillemotright}}1
}
\lstset{style=frode,basicstyle=\ttfamily\small}
\lstnewenvironment{kode}[1][]{\lstset{style=frode,#1}}{}
\newcommand{\cc}[1]{\texttt{#1}}            % kode i tabeller
\newcommand{\bs}{\textbackslash}

% ---------------------------------------------------------------- bokser ---
\tcbset{fbase/.style={enhanced,boxrule=0.6pt,arc=2pt,
  left=5pt,right=5pt,top=3pt,bottom=3pt,fonttitle=\sffamily\bfseries\small,
  before skip=6pt, after skip=6pt}}
\newtcolorbox{tips}[1][Tips]{fbase,colback=green!4,colframe=green!40!black,title={#1}}
\newtcolorbox{felle}[1][Felle]{fbase,colback=red!4,colframe=red!55!black,title={#1}}
\newtcolorbox{regel}[1][Regel]{fbase,colback=frodeblue!5,colframe=frodeblue,title={#1}}
\newtcolorbox{frode}[1][Frode sier]{fbase,colback=yellow!7,colframe=orange!70!black,title={#1}}

% ----------------------------------------------------------- overskrifter ---
\titleformat{\section}
  {\Large\sffamily\bfseries\color{frodeblue}}
  {\colorbox{frodeblue}{\color{white}\,\thesection\,}}{0.6em}{}
  [\vspace{-0.5em}{\color{frodeblue}\rule{\linewidth}{0.8pt}}]
\titleformat{\subsection}
  {\large\sffamily\bfseries\color{frodeblue!85!black}}{\thesubsection}{0.5em}{}
\titleformat{\subsubsection}
  {\normalsize\sffamily\bfseries}{\thesubsubsection}{0.5em}{}
\titlespacing*{\section}{0pt}{0pt}{10pt}
\titlespacing*{\subsection}{0pt}{12pt}{4pt}
\titlespacing*{\subsubsection}{0pt}{8pt}{2pt}
\newcommand{\kapittel}[1]{\clearpage\section{#1}}
% Bredere nummerfelt i innholdsfortegnelsen, så «12.10» ikke kolliderer med tittelen.
\makeatletter
\renewcommand*\l@subsection{\@dottedtocline{2}{1.5em}{2.8em}}
\makeatother

% Skilleside for Del 1 / Del 2 (egen side uten topptekst, med i innhold og bokmerker).
% Bruk: \del{Del 1 -- Oppgave 1}{Undertittel}
\newcommand{\del}[2]{%
  \clearpage\thispagestyle{empty}%
  \phantomsection\addcontentsline{toc}{part}{#1}%
  \vspace*{\fill}
  \begin{center}
    {\LARGE\sffamily\bfseries\color{frodeblue}#1\par}
    \vspace{8pt}\rule{0.45\linewidth}{0.6pt}\par\vspace{6pt}
    {\small #2\par}
  \end{center}
  \vspace*{\fill}\vspace*{3cm}
  \clearpage}

% Avkrysningsboks til sjekklistene (kapittel 1, 7 og 16).
\newcommand{\boks}{\raisebox{-0.2ex}{\setlength{\fboxsep}{0pt}\fbox{\rule{0pt}{1.5ex}\hspace{1.5ex}}}\ }

% ------------------------------------------------------------ topptekst ---
\pagestyle{fancy}
\fancyhf{}
\renewcommand{\sectionmark}[1]{\markboth{\thesection\quad #1}{}}
\fancyhead[L]{\sffamily\small\color{frodeblue}\textbf{Frode-bibelen}\ \textperiodcentered\ PROG1001}
\fancyhead[R]{\sffamily\small\leftmark}
\fancyfoot[C]{\sffamily\small\thepage}
\renewcommand{\headrulewidth}{0.4pt}

\setlist{nosep,leftmargin=1.4em}
\setlength{\parindent}{0pt}
\setlength{\parskip}{4pt}
\renewcommand{\arraystretch}{1.25}
\newcolumntype{L}[1]{>{\raggedright\arraybackslash}p{#1}}
